\documentclass[11pt]{article}
\usepackage[utf8]{inputenc}
\usepackage[margin=1in]{geometry}
\usepackage{hyperref}
\usepackage{setspace}

\begin{document}

\noindent
Carson Slater \hfill web: \href{https://carsonslater.github.io}{https://carsonslater.github.io} \hfill carsonslater7@gmail.com \\
\rule{\textwidth}{0.4pt}

\vspace{2em}

\begin{center}
    \Large \textbf{Cover Letter: Quantitative Analyst, Applied Biomechanics -- Philadelphia Phillies}
\end{center}

\vspace{1em}

\noindent
To the Hiring Manager, Applied Biomechanics Department,

\vspace{1em}

Baseball was my first love. From the time I was seven years old until my senior year of college, I spent countless hours in cleats, ultimately earning a spot in the starting lineup for Wheaton College's NCAA Division III varsity team. I eventually hung up my beloved baseball spikes to pursue a Ph.D. in Statistical Science at Baylor University, fully committing to the rigorous path of becoming a top-tier quantitative researcher. However, the opportunity to merge my domain knowledge of the game with my expertise in machine learning and statistical modeling as a Quantitative Analyst for the Philadelphia Phillies is a dream I couldn't pass up.

During my undergraduate studies, my transition into mathematics was largely fueled by exploring questions on the diamond. Whether it was building neural networks and linear regression models to investigate predictors for pitcher ERAs or utilizing ANOVA to analyze the impact of the universal designated hitter implementation, baseball data was the catalyst for my analytical journey. Later, as President of the Baylor University Sports Analytics Club, I continued to blend these passions by teaching version control, R programming, and machine learning to undergraduates while leading a team in the Kaggle March Madness bracket challenge.

In my current Ph.D. research at Baylor, I focus heavily on developing unsupervised learning methods for massive, messy time-series datasets. Specifically, my work funded by Denver Water involves engineering features from IoT meter data and building predictive models for hundreds of thousands of users. This experience translates seamlessly to handling high-frequency, complex biomechanical data in sports. I am highly proficient in building scalable research pipelines using R, Python, and SQL, and I am comfortable deploying these solutions utilizing ML best practices like Docker and Git.

The Quantitative Analyst role demands an individual who deeply understands both the nuances of statistical methodology---from regularization and ensemble methods to time series analysis---and the practical realities of player development and on-field outcomes. My background has equipped me to translate complex modeling into objective, actionable insights that can assist coaches and decision-makers in refining gameplanning and workload management.

Thank you for your time and consideration. I would welcome the opportunity to discuss how my technical background and lifelong passion for baseball can contribute to the Phillies' Applied Biomechanics team.

\vspace{1em}
\noindent
Sincerely,

\vspace{1em}
\noindent
\textbf{Carson Slater} \\
PhD Candidate, Statistical Science \\
Baylor University

\end{document}
